% =========================================================================
% SECCIÓN 5: BRECHAS DE INVESTIGACIÓN Y OPORTUNIDADES
% =========================================================================

\section{Brechas de Investigación y Oportunidades}
\label{sec:brechas_oportunidades}

\subsection{Brechas Críticas Identificadas}

\subsubsection{Brecha 1: Inexistencia de SLMs adaptados a variantes específicas del Quechua}
\noindent [Ausencia de modelos compactos afinados sobre Quechua Collao; la mayoría de trabajos agrupan variantes quechuas de forma homogénea, ignorando diferencias fonológicas y dialectales.]

\subsubsection{Brecha 2: Inadecuación de técnicas estándar de fragmentación RAG}
\noindent [Los algoritmos convencionales de chunking (por caracteres o párrafos) fragmentan sufijos aglutinantes y raíces semánticas clave, mermando la precisión de búsqueda vectorial.]

\subsubsection{Brecha 3: Carencia de pipelines RAG offline autosuficientes para entornos rurales}
\noindent [Predominio de arquitecturas dependientes de APIs comerciales externas (OpenAI, Claude), incompatibles con zonas rurales de baja o nula conectividad en Cusco y Puno.]

\subsubsection{Brecha 4: Falta de andamiaje pedagógico validado para EIB}
\noindent [Inexistencia de módulos de mediación pedagógica que guíen al estudiante en formato socrático bilingüe en lugar de suministrar transcripciones mecánicas.]

\subsubsection{Brecha 5: Vacío de benchmarks estandarizados y métricas culturales}
\noindent [Ausencia de suites de evaluación automatizadas (tipo MMLU o RAGAS) calibradas con criterios léxicos y culturales de la cosmovisión andina.]

\subsection{Oportunidades de Contribución para la Propuesta de Tesis}

\subsubsection{Oportunidad tecnológica: Arquitectura local híbrida SLM + RAG}
\noindent [Diseño de una arquitectura reproducible ejecutable en GPUs de consumo doméstico (4\,GB VRAM) y CPUs estándar.]

\subsubsection{Oportunidad lingüística: Curaduría y validación léxica Quechua Collao}
\noindent [Construcción de un banco curado de documentos educativos normalizados conforme a la ortografía oficial del MINEDU.]

\subsubsection{Oportunidad metodológica: Marco holístico de evaluación (Técnica + Pedagógica)}
\noindent [Formulación de un protocolo de validación que integre métricas computacionales (latencia, chrF++) con rúbricas de juicio de expertos en EIB.]
